\section{不可求导的函数}

这一节只是一个简单的介绍和一些可以争论的议论。

我们首先讨论两件似乎无关的事：第一件是回顾一下19世纪中分析数学发展的困难是怎样解决的；第二件是19世纪末，20世纪初物理学中的一件大事。

我们时常说19世纪是分析数学中注入了严格性的时代。这个时代的一位有重大贡献的代表性人物阿贝尔在1826年说过：“人们在分析中确实发现了惊人的含糊不清之处。这样一个完全没有计划和体系的分析，竟有那么多人能研究过它，真是奇怪。最坏的是，从来没有严格地对待过分析，非常奇怪的是这种方法只导致极少几个所谓悖论”（引文转引自克莱因《古今数学思想》中译本，第四卷，40章，上海科技出版社，1981年出版）。阿贝尔这样说是很自然的，即以函数概念为例，我们认识函数确实是由我们的感觉经验而非严格的公理化定义开始的。但是，这种感觉经验与大自然的规律符合到何种程度？它归根结蒂是不是也只是某种观测仪器给我们的经验曲线？人们总是不自主地想对“函数”加上种种限制。人们最容易接受的首先是连续性，连续性局部地看起来就是不间断和连通。由此自然想到，连续曲线下面的曲边梯形应该有面积 $\int_a^b f(x)\,dx$，所以人们以为只有可积的函数才是自然的，才是合理的，否则都是病态的。经验告诉我们曲线在每一点都有方向，例如质点运动的轨道总有方向。人们确实“感觉”到了 $df/dx$，同样人们也可以“感觉”到 $d^2f/dx^2$，因为一块石头或者一片树叶掉到头上造成的“感觉”确实是不一样的。于是人们又认为只有光滑函数才不是病态的。所以雅可比把这些光滑的函数称为“合理的函数”。但是人究竟能不能“感觉”到一个函数是 $C^4$ 而不是 $C^5$ 函数？人能不能依据“感觉”来区别 $C^\infty$ 函数与解析函数？实际上，人是越来越多地依靠逻辑，看由此形成的概念会得出什么样合理的结论，看是否能更深刻地揭示自然的规律。“不识庐山真面目，只缘身在此山中”，所以到19世纪的阿贝尔和今天的我们才会在瞠叹含糊不清之处多得令人吃惊的同时，也感到奇怪，这么多的含糊不清之处，竟然只形成“极少几个悖论”，竟然没有形成完全错误需要抛弃的理论。但是这种情况终于到了不能不解决的程度。我们下面连续三章都可以说是为了说明这一点。有些定理的叙述和证明将比前几章更细致，以避免出现漏洞。然而，为什么数学科学中“错误”一直较少呢？这不能不归功于在数学中有逻辑推理这样的工具，有严格性这样一种标准。回想一下，二千多年前的欧多克萨斯用不等式绕过了芝诺关于运动的悖论。现在人们发现，根本没有什么必要讨论无穷小，不必用玄学的语言来添乱，只要有极限理论一切都解决了。于是柯西说：$\lim_{x\to a}f(x)=A$ 就是“当 $x$ 越来越接近 $a$ 时，$f(x)$ 越来越接近 $A$”，而魏尔斯特拉斯还认为“越来越接近”的说法不清楚，仍然留下了一些来自运动学的概念的痕迹，而应该用 $\varepsilon-\delta$ 的语言，把问题归结为不等式。从这一点说，魏尔斯特拉斯是回到了欧多克萨斯，而把来自运动学的概念和语言诸如时间、运动等等，都“清”出了数学。从这一点讲，“变量”的数学变成了静态的数学。其实这是很自然的，因为例如时间更多的是物理学的研究对象。把时间引入数学就不可避免地带来了与此相关的物理学中的困难。所以到了19世纪末，人们就在这样一种十分讲究严格性的气氛中来讨论数学，其中自然也包括微积分。从表面上看，人们总在作“反面文章”：你说什么东西可以求导，数学家就找一个“反例”说这不一定行。人们想求极限，他又找一个“反例”，说求了极限就会出问题。所以许多定理甚至实际叙述也变得几乎看不懂，微积分变得更不近人情了，$\varepsilon-\delta$ 如噩梦似的压在学生们的心头。许多人作了种种努力使数学更“人性化”一些，但是没有一个人主张读牛顿、欧拉的“原汁原味”的微积分，而仍然要读魏尔斯特拉斯等人“消了毒”的微积分。

终于出现了康托尔（G. Cantor，1845--1918）的集合论。本书有一个界限，即不来讨论集合论的问题。我们只想指出，康托尔一生生活在争论之中，且集合论产生的“悖论”可不能说只有“极少几个”了，而且其影响深远。对于集合论，20世纪的两位数学巨人有截然不同的看法，一位是希尔伯特，他说康托尔的集合论为我们建立了数学的天堂乐园，现在谁也不能把我们从这个乐园中赶走。另一位同样伟大的数学家是庞加莱。他在《科学与方法》（1913）一书第二编“数学推理”中说了好长一段话：“逻辑有时造就奇异的东西。半个世纪以来，我们看到了出现一大堆异乎寻常的函数，它们似乎力图尽可能不类似于具有某些实效用的普通函数。它们不再有连续性，或者也许有连续性，但却没有导数等等。不仅如此，从逻辑的观点来看，正是这些奇异函数才是最普遍的。……从前，当一个新函数被发明时，正是为了实用的目的，今天为了指出我们祖先推理的错误，才特意发明新函数，人们从来也不能从中得到比这更多的东西。”说白了，这就是说现在人们研究数学为的只是证实他们的爷爷无能而且不严密。总之，数学是否过度强调严格化、公理化……妄想“清洗”直观的直觉，以至于远离对大自然的深刻的探索呢？这样，我们来看第二件事。

如果今天还有人不承认分子和原子的存在，那就是很可笑的事。但是当1897年元旦奥地利物理学家和哲学家马赫（E. Mach，1838--1916）（此人在中国名声不佳）在维也纳科学院一次会议上宣布“我不相信原子存在”时，得到的却不是嘲笑，而是有人反对，有人深思。这是物理学一次大辩论中的一幕。这场大争论终于导致物理学的大发展。问题在于，尽管当时从原子“假说”中已得出了许多正确的、与实验相符的结论，原子究竟还是不可捉摸的东西。对于原子是什么，有什么性质，还一直没有“直接”的证据。马赫从他的实证哲学的观点看来，原子的存在还只是一种形而上学的假说，所以应该抛弃。“直接”的证据出现在布朗运动的研究。布朗（Robert Brown）是英国植物学家，他注意到花粉在水面上的无规则运动（1827）。对此，只能用分子的碰撞来解释。爱因斯坦（1905）做了许多工作，并且计算出诸如阿伏伽德罗常数和分子的平均自由程，后来又经著名的法国实验物理学家佩兰（Jean Baptiste Perrin）系统的实验证实，原子是否存在的争论，才算尘埃落定。正是这位佩兰在他的名著《原子》一书中说：“利用显微镜观察布朗运动，看到浮在溶液中的每一个小颗粒来回运动时，我们被实验的真实完全吸引住了。为了对它的轨迹作切线，需要作连接轨迹上非常接近的两点（即该小颗粒在两邻近时刻的位置）的直线，并找出它的极限位置。可是在我们研究过程中，两次观察的时间间隔无论多小，这直线的方向总是持续地变化着。由这项研究，一个不带成见的观察者只能得到函数没有导数的感觉，而终究得不到有切线的曲线的感觉。”这段话转引自克莱因的名著《高观点下的初等数学》第三卷，湖北教育出版社，1993年，中译本46页。译者是已故的著名数学家吴大任教授和夫人陈骊教授。同页上还引述了博雷尔（E. Borel——测度理论的首创者之一）的一篇评论：“分子理论及其数学”，可惜无法找到原文。布朗运动的重要性早就不限于研究花粉的随机活动，而在随机现象的数学与物理理论乃至金融理论中都极为重要，而几乎所有布朗运动的轨道又都是处处不可求导的函数。当然这不是说不可求导函数的研究都是这些“实际需要”的“期货”。这样说，不仅不真实，而且不严肃。到了20世纪70年代末、80年代初，数学中出现了所谓分形几何学的研究，指出处处没有切线的曲线等等，在大自然中是十分普遍而且重要的。分形几何学在数学中也有长久的源流，与许多数学分支有关，所以我们没有力量去涉及有关的问题，而只想指出，正是由庞加莱开创的天体力学的研究中，在研究天体运动轨道的渐近状况时，就会出现构造极为奇特的集合。分形理论与它有密切关系。可见，数学走到了这一步，这种严格化、公理化的倾向其实标志了对大自然探索的深化。正是庞加莱自己的研究不但不是“为了指出我们祖先推理的错误”，反而证明了这种批判性的态度的远见。更令人吃惊的是布朗运动的轨道确与魏尔斯特拉斯的构造的不可求导的连续函数颇为相近！

魏尔斯特拉斯构造的函数是
\begin{equation}
 f(x)=\sum_{n=0}^{\infty}b^n\cos(a^n\pi x),
 \qquad -\infty<x<\infty,\quad 0<b<1.
 \tag{1}
\end{equation}
这里 $a>0$，其值待定。级数(1)显然在 $-\infty<x<+\infty$ 中一致收敛，所以 $f(x)$ 在 $-\infty<x<+\infty$ 上连续。它是无穷多个正弦曲线（余弦函数也是正弦曲线）叠加而成的，记其第 $n+1$ 条曲线为
\[
 y=T_n(x)=b^n\cos(a^n\pi x),
\]
其周期为 $2/a^n$，振幅为 $b^n$。我们知道，正弦曲线之斜率（按绝对值计）之最大值出现在它的零点处。我们将用这个最大值刻画它陡峭的程度，故不妨称之为 $T_n(x)$ 的陡度，于是 $T_n(x)$ 的陡度就是
\[
 \max\left|\frac{dT_n(x)}{dx}\right|=(ab)^n\pi.
\]
现在令
\begin{equation}
 a>1\text{ 为一整数，且 }ab>1,
 \tag{2}
\end{equation}
所以 $T_{n+1}(x)$ 之陡度又比 $T_n(x)$ 的陡度大 $ab$ 倍，所以构成函数(1)的这些正弦波（即正弦曲线）就越来越窄，越来越陡，其振幅也越来越小。图3-7-1就 $b=\frac12,a=5$ 画出了级数(1)的三个部分和。短划线是部分和 $S_0(x)=T_0(x)=\cos\pi x$；虚线是部分和 $S_1(x)=T_0(x)+T_1(x)=\cos\pi x+\frac12\cos5\pi x$；实线则是部分和 $S_2(x)=T_0(x)+T_1(x)+T_2(x)=\cos\pi x+\frac12\cos5\pi x+\frac14\cos25\pi x$。图3-7-1给我们的启发是：$T_n(x)$ 的零点与 $T_{n+1}(x)$ 的零点纠合在一起，这就是为什么我们要求 $a$ 为整数的原因。我们还要求 $T_n(x)$ 的极值点（即 $dT_n(x)/dx$ 的零点）又与 $T_{n+1}(x)$ 的极值点纠合在一起。因为 $T_n(x)$ 的每个周期内必有 $T_{n+1}(x)$ 的 $a$ 个周期，前者的每个半周中自然含有 $T_{n+1}(x)$ 的 $a$ 个半周。但是 $T_n(x)$ 的最大与最小值交替出现，两个相邻最大与最小值又只相距半周。所以若 $T_n(x)$ 与 $T_{n+1}(x)$ 有一个共同的最大值 $x$，则当 $x$ 从 $T_n(x)$ 的这个最大值走到下一个极值为最小值时，则 $T_{n+1}$ 也由最大变为最小，再变为最大……一共变了 $a$ 次。如果 $a$ 是奇数，则 $x$ 对于 $T_{n+1}(x)$ 而言仍然是由极大变成了极小。正是这样才能保证 $T_n(x)$ 叠加 $T_{n+1}(x)$ 时是极大叠加到极大上，极小叠加到极小上而不至于互相抵消。我们可以想像，叠加的结果会成为一条“毛茸茸”的曲线，而有可能是一个不可求导的函数的图像。所以我们除(1)以外还假设
\begin{equation}
 a>1\text{ 为一奇数}.
 \tag{3}
\end{equation}

\begin{figure}[htbp]
\centering
\begin{tikzpicture}[x=4.3cm,y=2.3cm,bella figure]
  \draw[->] (-1.05,0)--(1.10,0) node[right] {$x$};
  \draw[->] (0,-1.85)--(0,1.85);
  \draw[dashed,thick,domain=-1:1,samples=180,smooth,variable=\x]
     plot ({\x},{cos(180*\x)});
  \draw[dotted,thick,domain=-1:1,samples=260,smooth,variable=\x]
     plot ({\x},{cos(180*\x)+0.5*cos(900*\x)});
  \draw[thick,domain=-1:1,samples=700,smooth,variable=\x]
     plot ({\x},{cos(180*\x)+0.5*cos(900*\x)+0.25*cos(4500*\x)});
  \begin{scope}[shift={(0.38,1.52)},x=1cm,y=1cm]
    \draw[dashed,thick] (0,0)--(.62,0); \node[anchor=west] at (.70,0) {$y_0=\cos\pi x$};
    \draw[dotted,thick] (0,-.22)--(.62,-.22); \node[anchor=west] at (.70,-.22) {$y_1=\cos\pi x+\frac12\cos5\pi x$};
    \draw[thick] (0,-.44)--(.62,-.44); \node[anchor=west] at (.70,-.44) {$y_2=\cos\pi x+\frac12\cos5\pi x+\frac14\cos25\pi x$};
  \end{scope}
\end{tikzpicture}
\caption*{图3-7-1}
\end{figure}

下面把这些讨论“数学地”讲明白。

首先看 $T_n(x)$ 的零点
\[
 x_k^n=\frac{2k+1}{2a^n},
\]
其中 $k$ 是整数。这时 $a^{n+1}x_k^n=\frac{2k+1}{2}a$ 也是一个奇数的一半，所以
\[
 T_{n+1}(x_k^n)=b^{n+1}\cos\left(a^{n+1}x_k^n\pi\right)
 =b^{n+1}\cos\left(\frac{2k+1}{2}a\pi\right)=0,
\]
所以 $T_n(x)$ 的零点一定是 $T_{n+1}(x),T_{n+2}(x),\cdots$ 的零点，但不一定是 $T_0,\cdots,T_{n-1}$ 的零点。在函数(1)的图像中 $T_n(x)$ 的零点总是把 $T_n(x),T_{n+1}(x),\cdots$ 的零点纠合在一起，所以我们称 $x_k^n$ 为 $f(x)$ 之结点，并以
\[
 N=\left\{\frac{2k+1}{2a^n},\ k=0,\pm1,\pm2,\cdots,\ n=0,1,2,\cdots\right\}
\]
记结点之集合，我们看到

\textbf{引理1}\quad $N$ 是 $\mathbb R$ 的一个可数的稠密子集。

再看 $T_n(x)$ 之极值点 $\tilde x_k^n=k/a^n$，$k$ 是整数，这时 $a^{n+1}\tilde x_k^n=ak$ 故
\[
 T_n(\tilde x_k^n)=b^n\cos k\pi=(-1)^k b^n
\]
是 $T_n(x)$ 之极大（小）值，而
\[
 T_{n+1}(\tilde x_k^n)=b^{n+1}\cos ak\pi=(-1)^{ak}b^{n+1}.
\]
故 $\tilde x_k^n$ 也使 $T_{n+1}(x)$ 达到极大（小），而且与 $T_n(\tilde x_k^n)$ 有相同符号，所以极值点也是纠结在一起的，不会互相抵消；又因为各有不同因子 $b^n$，$0<b<1$，求和以后也不会变为正负无穷。我们称 $\tilde x_k^n$ 为 $f(x)$ 的峰点，并以
\[
 M=\left\{\frac{k}{a^n},\ k=0,\pm1,\pm2,\cdots,\ n=0,1,2,\cdots\right\}
\]
记峰点之集合，也很容易看到

\textbf{引理2}\quad $M$ 是 $\mathbb R$ 的一个可数稠密子集。

现在来证明函数(1)处处没有导数。为此任取一点 $x_0$ 并记(1)在 $x_0$ 之值为 $y_0=f(x_0)$，并且考虑 $T_n(x)$。$x_0$ 必在 $T_n(x)$ 的两个峰点之间，即有一个整数 $\alpha_n$ 使
\[
 \frac{\alpha_n}{a^n}<x_0<\frac{\alpha_n+1}{a^n}.
\]
在这两个峰点中取更接近 $x_0$ 的一个为 $\alpha_n/a^n$（若 $x_0$ 是中点，则取左方的一个），则有
\begin{equation}
 -\frac1{2a^n}<x_0-\frac{\alpha_n}{a^n}\le\frac1{2a^n},
 \qquad
 -\frac12<a^n x_0-\alpha_n\le\frac12.
 \tag{4}
\end{equation}
$\alpha_n/a^n$ 也有两个相邻结点
\[
 x'=\frac{\alpha_n-1}{a^n},
 \qquad
 x''=\frac{\alpha_n+1}{a^n}.
\]
令函数在这两点上之值为 $f(x')=y'$, $f(x'')=y''$，于是得到两个差商
\begin{equation}
 \frac{y'-f(x_0)}{x'-x_0},
 \qquad
 \frac{y''-f(x_0)}{x''-x_0}.
 \tag{5}
\end{equation}
为了证明 $f(x)$ 在 $x_0$ 没有导数，只需证明当 $x',x''\to x_0$ 时上述两个差商没有相同极限，也不趋向同号的无穷大。具体说来，我们要证明

\textbf{引理3}\quad 若 $0<b<1$，$a$ 为大于1的奇数且
\begin{equation}
 ab>1+\frac32\pi,
 \tag{6}
\end{equation}
则(5)中的两个差商不可能有相同的极限，也不趋向同号的无穷大。

\textbf{证}\quad 记 $x_{n+1}=x_0a^n-\alpha_n$，由(4)有 $|x_{n+1}|\le\frac12$，而且
\[
 x'-x_0=\frac{\alpha_n-1-x_0a^n}{a^n}
 =\frac1{a^n}[-1-(x_0a^n-\alpha_n)]
 =\frac1{a^n}(-1-x_{n+1}),
\]
\[
 x''-x_0=\frac{\alpha_n+1-x_0a^n}{a^n}
 =\frac1{a^n}[1-(x_0a^n-\alpha_n)]
 =\frac1{a^n}(1-x_{n+1}).
\]
因为 $a>1$，故当 $n$ 充分大时，$|x'-x_0|$ 与 $|x''-x_0|$ 皆可任意小。

再注意到
\[
 f(x)=\sum_{m=0}^{n-1}b^m\cos(a^m\pi x)
      +\sum_{m=n}^{\infty}b^m\cos(a^m\pi x)
      =S_{n-1}(x)+R_n(x).
\]
所以
\begin{align*}
 f(x')
 &=\sum_{m=0}^{n-1}b^m\cos(a^m\pi x')
   +\sum_{m=n}^{\infty}b^m\cos[(\alpha_n-1)a^{m-n}\pi]\\
 &=\sum_{m=0}^{n-1}b^m\cos(a^m\pi x')
   -(-1)^{\alpha_n}b^n\sum_{\nu=0}^{\infty}b^\nu,
\end{align*}
\begin{align*}
 f(x'')
 &=\sum_{m=0}^{n-1}b^m\cos(a^m\pi x'')
   +\sum_{m=n}^{\infty}b^m\cos[(\alpha_n+1)a^{m-n}\pi]\\
 &=\sum_{m=0}^{n-1}b^m\cos(a^m\pi x'')
   -(-1)^{\alpha_n}b^n\sum_{\nu=0}^{\infty}b^\nu.
\end{align*}
注意我们这里利用了 $a^{m-n}$ 为奇数 $2l+1$，故
\[
 \cos[(\alpha_n\pm1)a^{m-n}\pi]
 =\cos[(\alpha_n\pm1)\pi]
 =(-1)^{\alpha_n\pm1}
 =-(-1)^{\alpha_n}
\]
（$\alpha_n$ 是整数）。同样 $f(x_0)$ 也可以分成两项：利用 $x_{n+1}=x_0a^n-\alpha_n$，而 $x_0=a^{-n}(x_{n+1}+\alpha_n)$，故
\begin{align*}
 f(x_0)
 &=\sum_{m=0}^{n-1}b^m\cos(a^m\pi x_0)
   +\sum_{m=n}^{\infty}b^m\cos[a^{m-n}\pi(x_{n+1}+\alpha_n)]\\
 &=\sum_{m=0}^{n-1}b^m\cos(a^m\pi x_0)
   +(-1)^{\alpha_n}b^n\sum_{\nu=0}^{\infty}b^\nu\cos(a^\nu\pi x_{n+1}).
\end{align*}
于是(5)中的两个差商可以写成
\begin{align*}
 \frac{y'-f(x_0)}{x'-x_0}
 &=\sum_{m=0}^{n-1}b^m
   \frac{\cos(a^m\pi x')-\cos(a^m\pi x_0)}{x'-x_0}\\
 &\quad+b^n\sum_{\nu=0}^{\infty}b^\nu(-1)^{\alpha_n+1}
 \frac{1+\cos(a^\nu\pi x_{n+1})}{x'-x_0},
\end{align*}
\begin{align*}
 \frac{y''-f(x_0)}{x''-x_0}
 &=\sum_{m=0}^{n-1}b^m
   \frac{\cos(a^m\pi x'')-\cos(a^m\pi x_0)}{x''-x_0}\\
 &\quad+b^n\sum_{\nu=0}^{\infty}b^\nu(-1)^{\alpha_n+1}
 \frac{1-\cos(a^\nu\pi x_{n+1})}{x''-x_0}.
\end{align*}
现在分别估计各个差商的两项。对每一个差商的第一项和式利用和差化积公式有
\begin{align*}
 &\sum_{m=0}^{n-1}b^m
 \frac{\cos(a^m\pi x')-\cos(a^m\pi x_0)}{x'-x_0}\\
 &\qquad=\sum_{m=0}^{n-1}\frac{b^m}{x'-x_0}
 \left[-2\sin\!\left(a^m\pi\frac{x'+x_0}{2}\right)
 \sin\!\left(a^m\pi\frac{x_0-x'}{2}\right)\right]\\
 &\qquad=\sum_{m=0}^{n-1}-a^m b^m\pi\,
 \frac{\sin\!\left(a^m\pi\frac{x_0-x'}2\right)}
      {a^m\pi\frac{x_0-x'}2}
 \sin\!\left(a^m\pi\frac{x'+x_0}{2}\right).
\end{align*}
因为
\[
 \left|\frac{\sin\!\left(a^m\pi\frac{x_0-x'}2\right)}
 {a^m\pi\frac{x_0-x'}2}\right|\le1,
 \qquad
 \left|\sin\!\left(a^m\pi\frac{x'+x_0}{2}\right)\right|\le1,
\]
考虑到 $a^m b^m\pi>0$，所以这一和式的绝对值
\[
 \le \pi\sum_{m=0}^{n-1}a^m b^m
 =\pi\frac{a^n b^n-1}{ab-1}
 <\frac{\pi a^n b^n}{ab-1}.
\]
因此存在一个常数 $\varepsilon$，$-1<\varepsilon<1$，使
\begin{equation}
 \text{第一个和式}=\frac{\varepsilon\pi a^n b^n}{ab-1}.
 \tag{7}
\end{equation}
第二个差商的第一个和式是一个收敛的无穷级数。因为
\[
 x'-x_0=-\frac1{a^n}(1+x_{n+1}),
\]
所以
\begin{align*}
 &b^n\sum_{\nu=0}^{\infty}b^\nu(-1)^{\alpha_n+1}
 \frac{1+\cos(a^\nu\pi x_{n+1})}{x'-x_0}\\
 &\qquad=(-1)^{\alpha_n}a^n b^n
 \sum_{\nu=0}^{\infty}b^\nu
 \frac{1+\cos(a^\nu\pi x_{n+1})}{1+x_{n+1}}.
\end{align*}
现在新级数第一项即 $\nu=0$ 的一项是
\[
 \frac{1+\cos(\pi x_{n+1})}{1+x_{n+1}}>\frac23.
\]
这是因为 $|x_{n+1}|\le\frac12$，从而分子 $1+\cos(\pi x_{n+1})\ge1$，而分母 $\le\frac32$。总之新级数第一项 $\ge\frac23$，其余各项均为非负，所以其和 $\ge\frac23$。记其和为 $\eta'\frac23$，应有 $\eta'\ge1$。合并以上的结果，
\[
 \text{第一个差商}=\frac{\varepsilon\pi a^n b^n}{ab-1}
 +(-1)^{\alpha_n}\frac23\eta' a^n b^n.
\]
如果把 $\varepsilon$ 写成 $\varepsilon=(-1)^{\alpha_n}\eta'\varepsilon'$，则 $\varepsilon'=(-1)^{\alpha_n}\varepsilon/\eta'$，$|\varepsilon'|=|\varepsilon|/\eta'<1$，而
\begin{equation}
 \text{第一个差商}=(-1)^{\alpha_n}a^n b^n
 \left(\frac23+\frac{\varepsilon'\pi}{ab-1}\right)\eta'
 =(-1)^{\alpha_n}a^n b^n\rho'.
 \tag{8}
\end{equation}
这里
\[
 \rho'=\left(\frac23+\frac{\varepsilon'\pi}{ab-1}\right)\eta'
 >\frac23-\frac{\pi}{ab-1}>0.
\]
这里我们应用了(6)式。

第二个差商可以类似地处理，注意到这时
\[
 x''-x_0=\frac1{a^n}(1-x_{n+1}),
\]
我们将得到
\begin{equation}
 \text{第二个差商}=(-1)^{\alpha_n+1}a^n b^n
 \left(\frac23+\frac{\varepsilon''\pi}{ab-1}\right)\eta''
 =(-1)^{\alpha_n+1}a^n b^n\rho'',
 \tag{9}
\end{equation}
这里
\[
 \rho''=\left(\frac23+\frac{\varepsilon''\pi}{ab-1}\right)\eta''
 >\frac23-\frac{\pi}{ab-1}>0.
\]
总结以上即知：若 $(-1)^{\alpha_n}=1$，有
\[
 \frac{y'-f(x_0)}{x'-x_0}\ge a^n b^n\left(\frac23-\frac{\pi}{ab-1}\right),
\]
\[
 \frac{y''-f(x_0)}{x''-x_0}\le -a^n b^n\left(\frac23-\frac{\pi}{ab-1}\right),
\]
当 $(-1)^{\alpha_n}=-1$ 时
\[
 \frac{y'-f(x_0)}{x'-x_0}\le -a^n b^n\left(\frac23-\frac{\pi}{ab-1}\right),
\]
\[
 \frac{y''-f(x_0)}{x''-x_0}\ge a^n b^n\left(\frac23-\frac{\pi}{ab-1}\right).
\]
由于 $x',x''$ 属于 $\mathbb R$ 上的一个可数稠密集 $M$，故如函数(1)的差商不可能有极限，也不可能趋向于同符号的无穷大。魏尔斯特拉斯的函数(1)在任一点 $x_0$ 都没有导数得证。
